\section{从描述走向干预}

前几章已经论证，行为失败往往是一个状态转变问题，而不只是欲望不足。只有当这一主张能够改变设计实践时，它才真正重要。因此，干预章节承担的任务要比动机章节更严格：它必须明确说明改变的是哪一个控制变量，为什么这一变量比直接劝勉更便宜或更可靠，以及设计中的哪些部分得到了较强证据支持，而不只是停留在直觉层面。

在本书中，干预设计由五个观念支配：
\begin{enumerate}[nosep]
  \item 干预作用于转变，而不是作用于抽象的身份宣称；
  \item 决策窗口之所以重要，是因为杠杆对时机高度敏感；
  \item 重复暴露与历史依赖之所以重要，是因为框架不会在每次尝试时都从零重建；
  \item 吸引子与亚稳态语言有助于解释，为什么一些局部重设计会产生不成比例的下游影响；
  \item 来源纪律之所以重要，是因为经验发现、理论机制构建与推测性的前沿主张，不应被混同为具有同等地位。
\end{enumerate}